\documentclass{article}

\usepackage[utf8]{inputenc}
\usepackage[T1]{fontenc}
\usepackage{amsmath,amssymb,amsthm}
\usepackage{mathtools}
\usepackage{hyperref}
\usepackage{booktabs}

\makeatletter
\newcommand{\citep}[2][]{%
  \begingroup
  \def\@tmpkey{#2}%
  \ifx\@tmpkey\@empty\else
    (\ifx\@empty#1\@empty\else #1, \fi
     \@nameuse{bibkey@\@tmpkey})%
  \fi
  \endgroup
}
\expandafter\def\csname bibkey@lasser2026horizon\endcsname{Horizon~Aware}
\expandafter\def\csname bibkey@lasser2026exo\endcsname{Exogenous~Verification}
\expandafter\def\csname bibkey@lasser2026phase\endcsname{Phase~Redundancy}
\expandafter\def\csname bibkey@lasser2026paper8a\endcsname{Revealed~Sacrifice}
\expandafter\def\csname bibkey@lasser2026paper8b\endcsname{Need~Sufficiency}
\expandafter\def\csname bibkey@lasser2026micro\endcsname{Microfoundation}
\makeatother

\newtheorem{definition}{Definition}
\newtheorem{proposition}{Proposition}
\newtheorem{lemma}{Lemma}
\newtheorem{theorem}{Theorem}
\newtheorem{remark}{Remark}
\newtheorem{assumption}{Assumption}

\DeclareMathOperator{\vol}{vol}
\newcommand{\volL}{{\vol_{\mathrm{P}}}}
\newcommand{\volR}{{\vol_{\mathrm{R}}}}
\newcommand{\volRwin}{{\vol_{\mathrm{R}}^{[W]}}}
\newcommand{\volRlower}{{\vol_{\mathrm{R}}^{\mathrm{lower}}}}
\newcommand{\rext}{r_{\mathrm{ext}}}
\newcommand{\rsub}{r_{\mathrm{sub}}}
\newcommand{\rS}{r_{\mathrm{S}}}
\newcommand{\rW}{r_{\mathrm{W}}}
\newcommand{\rhomincross}{\rho_{\min}^{\mathrm{cross}}}
\newcommand{\betaL}{\beta^{\mathrm{lower}}}
\newcommand{\HHI}{\mathrm{HHI}}
\newcommand{\Pproxy}{P}
\newcommand{\Ttruth}{T}
\newcommand{\epsgap}{\varepsilon_{\mathrm{gap}}}
\newcommand{\epsnonres}{\varepsilon_{\mathrm{gap}}^{\mathrm{nonres}}}
\newcommand{\epsfloor}{\varepsilon_{\mathrm{floor}}^{\mathrm{res}}}
\newcommand{\epsnonresComposed}{\varepsilon_{\mathrm{gap},\mathrm{composed}}^{\mathrm{nonres}}}
\newcommand{\epsnonresChannel}[1]{\varepsilon_{\mathrm{gap},(#1)}^{\mathrm{nonres}}}
\newcommand{\epsnonresCoev}{\varepsilon_{\mathrm{gap},\mathrm{coev}}^{\mathrm{nonres}}}

\title{Lemma 1 (Intensive composition under co-evolution): Full Proof Draft}
\author{Paper 10 working draft for codex review}
\date{}

\begin{document}

\maketitle

\section{Goal of this draft}

Promote Lemma~1 of Paper~10 \S4.2 from inline proof outline to a
self-contained formal proof, suitable for placement in the
appendix. The current outline references source-paper machinery but
does not formalize the ``intensive in $|P|$'' property nor
explicitly verify that the composition rule preserves it through
the co-evolution error terms.

This draft is for codex review before integration into Paper~10's
appendix.

\section{Formal definition of intensivity}

\begin{definition}[Intensive in $|P|$]
\label{def:intensive}
A bound $B$ depending on the deployment state is \emph{intensive in
$|P|$} if there exists a constant $C \geq 0$, independent of the
capability poset size $|P|$, such that $B \leq C$ for all valid
deployment states.

A bound is \emph{extensive in $|P|$} if no such $|P|$-independent
constant exists; equivalently, $B$ admits states where $B \to
\infty$ as $|P| \to \infty$.
\end{definition}

\paragraph{Note.} Definition~\ref{def:intensive} treats intensivity
uniformly: the bounding constant $C$ may depend on the deployment's
threat model, training discipline, or other operational parameters,
but must not depend on $|P|$. This is the property the deployment
claim's capability-magnitude-independence requires.

\section{Source-paper intensivity assumptions}

We assume the following intensivity properties of the source-paper
quantities used in Paper~10's composition. Each is established in
its source paper or follows from the source-paper machinery; we
state them as named assumptions for clarity.

\begin{assumption}[A1: Per-channel non-residual gap intensivity]
\label{ass:A1}
For each channel $c \in \{1, 2, 3, 4\}$ in
\citep[Microfoundation]{lasser2026micro}'s four-channel decomposition,
the per-channel non-residual gap contribution $\epsnonresChannel{c}$
is intensive in $|P|$: there exists $K_c \geq 0$ independent of
$|P|$ such that $\epsnonresChannel{c} \leq K_c$ for all valid
deployment states.
\end{assumption}

\textbf{Justification.} \citep[Microfoundation]{lasser2026micro}
defines each channel's contribution per capability, with bounds
derived from per-capability quantities (observation density,
attestation quality, etc.). Each is bounded by the channel's
operational threshold, which depends on the deployment's
verification infrastructure (Paper 5's witness density, attestation
protocols), not on $|P|$.

\begin{assumption}[A2: Residual floor intensivity]
\label{ass:A2}
The residual floor $\epsfloor$ from
\citep[Microfoundation]{lasser2026micro}'s gap decomposition is
intensive in $|P|$: there exists $K_{\mathrm{floor}} \geq 0$
independent of $|P|$ such that $\epsfloor \leq K_{\mathrm{floor}}$
for all valid deployment states.
\end{assumption}

\textbf{Justification.}
\citep[Microfoundation]{lasser2026micro}'s residual class is the
structurally-unobservable subset of capability-space; its floor is
bounded by the operational unobservability rate (a per-capability
quantity bounded by the deployment's privacy and attestation
discipline), not by the absolute capability count.

\begin{assumption}[A3: Bounded Lipschitz alignment property]
\label{ass:A3}
The alignment property $g$ has Lipschitz constant
$\mathrm{Lip}(g)$ that is bounded by some $K_{\mathrm{Lip}} \geq 0$
independent of $|P|$. Equivalently, $g$ is in the class of
welfare functionals whose sensitivity to single-capability
perturbations is uniformly bounded.
\end{assumption}

\textbf{Justification.} The intensivity of the deployment claim
requires a class restriction on $g$. Standard welfare functionals
(e.g., scalarizations of per-capability metrics, weighted sums with
bounded weights) satisfy A3 by construction. Pathological
functionals where $\mathrm{Lip}(g)$ grows with $|P|$ would produce
extensive Goodhart bounds; A3 excludes such functionals from the
deployment claim's scope. In Paper~10, A3 is promoted to explicit
condition (C6) of Theorem~\ref{thm:deployment_safety}.

\begin{assumption}[A4: Co-evolution error bound]
\label{ass:A4}
By
\citep[Microfoundation]{lasser2026micro}'s Composition Proposition 1,
the co-evolution positive-part error term $(\epsnonresCoev)_+$ admits
the bound:
\[
  (\epsnonresCoev)_+ \;\leq\; K_{\mathrm{coev}} \cdot M(\text{deployment}),
\]
where $K_{\mathrm{coev}}$ is a structural constant from
\citep[Microfoundation]{lasser2026micro}'s proposition and
$M(\text{deployment})$ is the maximum per-channel coupling
magnitude (a per-channel rate quantity bounded by the deployment's
operational discipline, not by $|P|$).
\end{assumption}

\textbf{Justification.}
\citep[Microfoundation]{lasser2026micro}'s Composition Proposition 1
establishes that the co-evolution regime introduces approximate
composition with positive-part error terms whose magnitude depends
on inter-channel coupling rates. Coupling rates are per-channel
quantities (how rapidly Channel-3 individuation events affect
Channel-2 attestation, etc.), bounded by the deployment's
governance and discipline structure. They do not scale with $|P|$.

\section{Statement of Lemma 1}

\begin{lemma}[Intensive composition under co-evolution]
\label{lem:1}
Let the composition rule $B$ be the static Goodhart bound:
\begin{equation}
\label{eq:composition_rule}
  B \;=\; \mathrm{Lip}(g) \cdot \bigl(\epsnonresComposed + \lambda \cdot \epsfloor\bigr),
\end{equation}
where $\epsnonresComposed$ is the composed non-residual gap under
co-evolution, $\epsfloor$ is the residual floor, $\lambda \in [0,1]$
is the residual weight, and $\mathrm{Lip}(g)$ is the Lipschitz
constant of the alignment property $g$.

Under Assumptions~\ref{ass:A1}--\ref{ass:A4}, $B$ is intensive in
$|P|$: there exists a constant $C^* \geq 0$ independent of $|P|$
such that $B \leq C^*$ for all valid deployment states.
\end{lemma}

\section{Proof}

\begin{proof}
We compose Assumptions~\ref{ass:A1}--\ref{ass:A4} into a uniform
bound on $B$.

\emph{Step 1: Bound the composed non-residual gap via four-channel
decomposition.}
By
\citep[Microfoundation]{lasser2026micro}'s Composition Proposition 1
under the co-evolution regime:
\begin{equation}
\label{eq:eps_decomp}
  \epsnonresComposed \;\leq\; \sum_{c \in \{1,2,3,4\}}
  \epsnonresChannel{c} + (\epsnonresCoev)_+.
\end{equation}
This is the explicit four-channel decomposition with positive-part
co-evolution error.

\emph{Step 2: Apply A1 to bound the channel sum.}
By Assumption~\ref{ass:A1}, each $\epsnonresChannel{c} \leq K_c$
with $K_c$ independent of $|P|$. Therefore:
\begin{equation}
\label{eq:channel_sum_bound}
  \sum_{c \in \{1,2,3,4\}} \epsnonresChannel{c}
  \;\leq\; \sum_{c \in \{1,2,3,4\}} K_c
  \;=:\; K_{\mathrm{ch}}.
\end{equation}
$K_{\mathrm{ch}}$ is finite and independent of $|P|$.

\emph{Step 3: Apply A4 to bound the co-evolution error.}
By Assumption~\ref{ass:A4}, $(\epsnonresCoev)_+ \leq K_{\mathrm{coev}}
\cdot M(\text{deployment})$ with both $K_{\mathrm{coev}}$ and
$M(\text{deployment})$ independent of $|P|$. Therefore:
\begin{equation}
\label{eq:coev_bound}
  (\epsnonresCoev)_+ \;\leq\; K_{\mathrm{coev}} \cdot M
  \;=:\; K_{\mathrm{coev}}',
\end{equation}
where we absorbed $M(\text{deployment})$ into the constant.

\emph{Step 4: Combine Steps 2 and 3 into a bound on $\epsnonresComposed$.}
Substituting Equations~\ref{eq:channel_sum_bound}
and~\ref{eq:coev_bound} into Equation~\ref{eq:eps_decomp}:
\begin{equation}
\label{eq:eps_total_bound}
  \epsnonresComposed \;\leq\; K_{\mathrm{ch}} + K_{\mathrm{coev}}'
  \;=:\; K_{\mathrm{nonres}}.
\end{equation}
$K_{\mathrm{nonres}}$ is the sum of $|P|$-independent constants and
is therefore $|P|$-independent.

\emph{Step 5: Apply A2 to bound the residual floor.}
By Assumption~\ref{ass:A2}, $\epsfloor \leq K_{\mathrm{floor}}$
with $K_{\mathrm{floor}}$ independent of $|P|$.

\emph{Step 6: Bound $\epsgap$ via the gap decomposition.}
The total gap is bounded by the non-residual plus the residual:
\begin{equation}
\label{eq:eps_gap_total}
  \epsgap \;\leq\; \epsnonresComposed + \lambda \cdot \epsfloor
  \;\leq\; K_{\mathrm{nonres}} + \lambda \cdot K_{\mathrm{floor}}
  \;\leq\; K_{\mathrm{nonres}} + K_{\mathrm{floor}}
\end{equation}
(using $\lambda \in [0,1]$ in the last inequality). Define
$K_{\mathrm{gap}} := K_{\mathrm{nonres}} + K_{\mathrm{floor}}$;
$K_{\mathrm{gap}}$ is independent of $|P|$.

\emph{Step 7: Apply A3 to bound the Lipschitz constant.}
By Assumption~\ref{ass:A3}, $\mathrm{Lip}(g) \leq K_{\mathrm{Lip}}$
with $K_{\mathrm{Lip}}$ independent of $|P|$.

\emph{Step 8: Compose into a uniform bound on $B$.}
Combining Equations~\ref{eq:composition_rule} with the bounds from
Steps 6 and 7:
\begin{equation}
\label{eq:final_bound}
  B \;=\; \mathrm{Lip}(g) \cdot \bigl(\epsnonresComposed + \lambda
  \cdot \epsfloor\bigr)
  \;\leq\; K_{\mathrm{Lip}} \cdot K_{\mathrm{gap}}
  \;=:\; C^*.
\end{equation}
$C^*$ is the product of $|P|$-independent constants and is therefore
$|P|$-independent.

By Definition~\ref{def:intensive}, $B$ is intensive in $|P|$. $\Box$
\end{proof}

\section{Discussion of the constants}

The proof produces an explicit constant $C^* = K_{\mathrm{Lip}}
\cdot (K_{\mathrm{ch}} + K_{\mathrm{coev}}' + K_{\mathrm{floor}})$,
each component of which depends only on operationally-determined
quantities:

\begin{itemize}
\item $K_{\mathrm{Lip}}$: bounded sensitivity of the alignment
  property $g$ (deployment-class restriction; (C6) of Theorem~1).
\item $K_c$ for $c \in \{1,2,3,4\}$: per-channel observation
  thresholds (set by the deployment's verification infrastructure).
\item $K_{\mathrm{coev}}$: structural constant from
  \citep[Microfoundation]{lasser2026micro}'s Composition Proposition 1.
\item $M$: maximum per-channel coupling magnitude (set by the
  deployment's governance discipline).
\item $K_{\mathrm{floor}}$: residual unobservability rate (set by
  the deployment's privacy and attestation discipline).
\end{itemize}

None of these scale with $|P|$. The constant $C^*$ is therefore a
deployment-time parameter, fully determined by operational
choices, and the deployment claim's capability-magnitude-
independence property follows.

\section{Connection to the deployment claim}

Lemma~\ref{lem:1}'s intensivity property is the foundational input
to Theorem~1's Layer~1 bound. Specifically:
\begin{itemize}
\item Layer~1 statement: $|g(\Ttruth) - g(\Pproxy)| \leq
  \mathrm{Lip}(g) \cdot h_{\mathrm{static}}(\theta)$.
\item By the Lipschitz transfer of
  \citep[Microfoundation]{lasser2026micro}, $|g(\Ttruth) -
  g(\Pproxy)| \leq \mathrm{Lip}(g) \cdot \epsgap$.
\item By Lemma~\ref{lem:1}, $\mathrm{Lip}(g) \cdot \epsgap \leq
  C^*$ uniformly.
\item Therefore $h_{\mathrm{static}}(\theta) \leq C^* /
  \mathrm{Lip}(g) = K_{\mathrm{gap}}$, intensive in $|P|$.
\end{itemize}

Layer~1's intensive bound is established through Lemma~\ref{lem:1}.
Subsequent lemmas (Lemma~2 Lyapunov-Goodhart bridge, Lemma~3 HHI
surrogate adequacy) refine the operational meaning of $\epsgap$
without changing its intensivity.

\section{What this proof does and does not establish}

\textbf{Establishes:}
\begin{itemize}
\item A formal definition of intensivity in $|P|$.
\item Four named assumptions (A1--A4) sufficient for the
  composition's intensivity, with each grounded in source-paper
  machinery.
\item An explicit constant $C^*$ in terms of operational
  parameters, with no $|P|$-dependence.
\end{itemize}

\textbf{Does not establish (deferred):}
\begin{itemize}
\item Tightness of the constant $C^*$. The proof gives a uniform
  upper bound, not the tightest possible bound. Deployments may
  achieve smaller bounds through tighter operational discipline;
  this is operational engineering, not theorem content.
\item Existence of deployment configurations satisfying A1--A4.
  The assumptions are conditions on the deployment; whether they
  hold for any particular deployment is verified by the
  deployment-tooling specification (\S8 of Paper~10).
\item The explicit form of the constants $K_c$, $K_{\mathrm{coev}}$,
  $M$, $K_{\mathrm{floor}}$, $K_{\mathrm{Lip}}$. These are determined
  by the deployment's operational choices and source-paper
  parameters; the proof asserts their existence and independence
  from $|P|$ but does not derive their numerical values.
\end{itemize}

\section{Specific verification questions for codex review}

\begin{enumerate}
\item Is Definition~\ref{def:intensive} (intensive in $|P|$) the
  right formalization? In particular, is the uniform-constant
  formulation appropriate, or should we use a different convention
  (e.g., $|P|$-asymptotic rather than $|P|$-uniform)?

\item Are Assumptions A1--A4 stated correctly? In particular:
  (a) Does A1's per-channel intensivity actually follow from
  \citep[Microfoundation]{lasser2026micro}, or does it require
  additional content not in the source paper?
  (b) Is A2 (residual-floor intensivity) defensible, or could the
  residual class grow with $|P|$ in some deployment regimes?
  (c) Is A4's bound on $(\epsnonresCoev)_+$ actually established
  by \citep[Microfoundation]{lasser2026micro}'s Composition
  Proposition~1, or does that proposition give a different
  (potentially weaker) bound?

\item Is the proof step-by-step rigorous? In particular,
  (a) Does the absorption $K_{\mathrm{coev}} \cdot M(\text{deployment})
  =: K_{\mathrm{coev}}'$ in Step 3 hide any $|P|$-dependence?
  (b) Is the use of $\lambda \in [0,1]$ in Step 6 legitimate, or
  could $\lambda$ depend on $|P|$ in some deployment regimes?

\item The proof asserts the constants exist but does not derive
  their numerical values. Is that acceptable for the deployment
  claim's purposes, or does Paper~10 need to specify computational
  procedures for each constant?

\item Is the Layer~1 connection (\S6 of this draft) correctly
  stated? Specifically, does $h_{\mathrm{static}}(\theta) \leq
  K_{\mathrm{gap}}$ correctly capture the role of
  Lemma~\ref{lem:1} in Theorem~1's Layer~1 proof?
\end{enumerate}

\end{document}
